\section{Evaluation}
\label{sec:evaluation}

This chapter evaluates the proposed staged validation pipeline with respect to the documentation quality dimensions defined in Section~2.3 (structural completeness, integrity, semantic adequacy, and cross-view consistency). The evaluation is designed so that the study setup and the analysis protocol remain valid independently of the final numerical outcomes; quantitative result sections are completed once Pass~1 and Pass~2 are executed on the evaluation corpus.

\subsection{Evaluation objectives}
\label{sec:eval_objectives}

The evaluation assesses whether the pipeline improves the detection and localization of architecture documentation issues in a way that is auditable and practically useful. Three objectives guide the evaluation:

\begin{enumerate}
    \item \textbf{Stage contribution:} quantify and explain what each stage adds (baseline $\rightarrow$ cosine resolver $\rightarrow$ preprocessing $\rightarrow$ Pass~1 $\rightarrow$ Pass~2).
    \item \textbf{Auditability:} verify that outputs remain evidence-backed and actionable, consistent with the auditability constraints in Section~5.4.7.
    \item \textbf{Industrial relevance:} assess whether detected issues align with reviewer expectations and whether the output format can support real review workflows.
\end{enumerate}

\subsection{Case study and dataset setup}
\label{sec:eval_dataset}

\subsubsection{Industrial document corpus}
\label{sec:eval_corpus}

The evaluation corpus consists of real architecture documentation authored using an arc42-derived company template and maintained in AsciiDoc (\texttt{.adoc}) format. Documents are maintained per software unit/subsystem (e.g., SWE-level documentation) and are processed in batch form, consistent with the workflow described in Chapter~4.

\paragraph{Corpus scope (to fill).}
\begin{itemize}
    \item Number of SWE documents: \textbf{[N]}
    \item Snapshot date / time window: \textbf{[DATE]}
    \item Document selection criterion: \textbf{[e.g., latest mainline versions]}
    \item Template version identifier: \textbf{[Template ID / commit hash]}
\end{itemize}

\subsubsection{Reference template}
\label{sec:eval_template}

All validation is performed against the company-specific arc42-derived reference template described in Section~2.2. Template guidance is used both as normative expectation for deterministic structure checks and as the reference intent for semantic adequacy checks.

\paragraph{Template source (to fill).}
\begin{itemize}
    \item Template artifact: \texttt{expected\_spec.json} (derived from company template AsciiDoc)
    \item Template revision: \textbf{[commit hash / version tag]}
\end{itemize}

\subsubsection{Preprocessing and input normalization}
\label{sec:eval_preprocessing}

Before evaluation, documents are parsed and normalized into the structured JSON representation described in Chapter~6. The representation captures section/heading hierarchy, content blocks, and associated template guidance per section/heading. The evaluation uses persisted intermediate artifacts (baseline reports, cosine reports, and preprocessed inputs) to ensure that stage-by-stage comparisons are reproducible.

\subsection{Metrics, criteria, and evaluation protocol}
\label{sec:eval_metrics}

This section defines the evaluation measures and how results are interpreted.

\subsubsection{Structural metrics (Stage~1 baseline)}
\label{sec:eval_structural_metrics}

Structural completeness is measured using deterministic signals from the baseline reports:
\begin{itemize}
    \item Missing required sections (count per document; list of missing sections)
    \item Missing required headings (count per section; list of headings)
    \item Hierarchy mismatches (if enabled; count and types)
    \item Extra sections/headings (count; treated as deviation/extension signals)
\end{itemize}

\subsubsection{Robustness metrics (Stage~2 cosine resolver)}
\label{sec:eval_robustness_metrics}

Robustness to naming variation is measured by:
\begin{itemize}
    \item Rescued mappings count: number of template sections aligned to document sections via similarity
    \item Reduction in missing-section findings after applying rescue mappings (before $\rightarrow$ after)
    \item Precision sanity checks: sampled manual verification of rescued mappings (sampling strategy in Section~\ref{sec:eval_qualitative})
\end{itemize}

Additionally, optional indicators from the cosine layer are reported:
\begin{itemize}
    \item high-confidence unfilled placeholder headings detected,
    \item near-verbatim template guidance copies detected,
    \item duplicate content candidates detected.
\end{itemize}

\subsubsection{Integrity and semantic adequacy metrics (Pass~1)}
\label{sec:eval_pass1_metrics}

Pass~1 evaluates integrity and semantic adequacy at section/heading granularity. The following measures are reported:
\begin{itemize}
    \item Integrity issue count per document and per section (e.g., placeholders, TODOs, empty headings)
    \item Semantic mismatch count per document and per section (guidance not satisfied, irrelevant content, boilerplate)
    \item Evidence coverage rate: percentage of findings that include a valid evidence snippet
    \item Section score distribution (0--5) across documents and sections
\end{itemize}

In addition to quantitative counts, a qualitative inspection assesses whether findings are clearly justified and human-verifiable.

\subsubsection{Cross-view consistency metrics (Pass~2)}
\label{sec:eval_pass2_metrics}

Pass~2 evaluates cross-view consistency through entity extraction and conservative rule checks. Measures include:
\begin{itemize}
    \item Entity coverage: counts of extracted entities by type (units/building blocks, interfaces, external systems, stakeholders)
    \item Consistency issue count by rule category:
    \begin{itemize}
        \item Building Block View $\leftrightarrow$ Runtime View (undefined entities referenced in runtime scenarios)
        \item Building Block View $\leftrightarrow$ Deployment View (missing or unclear mappings)
        \item Goals $\leftrightarrow$ Decisions (contradictions or missing alignment where checkable)
        \item Decisions $\leftrightarrow$ Risks/Technical Debt (unreflected trade-offs where explicit evidence exists)
    \end{itemize}
    \item Evidence coverage rate for consistency issues
\end{itemize}

\subsubsection{Stage contribution analysis (primary evaluation lens)}
\label{sec:eval_stage_contribution}

The primary evaluation method is a stage comparison:
\begin{enumerate}
    \item Stage~1 only (baseline)
    \item Stage~1 + Stage~2 (baseline + cosine resolver)
    \item Stage~1 + Stage~2 + Stage~3 (preprocessed audit readiness)
    \item Pass~1 (integrity + semantic adequacy)
    \item Pass~2 (cross-view consistency)
\end{enumerate}

For each stage set, the evaluation reports:
\begin{itemize}
    \item additional issues found,
    \item issues resolved or reduced (e.g., fewer false ``missing'' sections after rescue),
    \item changes in section scores (where applicable),
    \item representative qualitative examples illustrating the incremental benefit.
\end{itemize}

\subsection{Results on real documents}
\label{sec:eval_results}

This section reports quantitative results once the pipeline is executed on the corpus.

\subsubsection{Dataset summary}
\label{sec:eval_dataset_summary}

\paragraph{To fill after execution.}
\begin{itemize}
    \item Documents processed successfully: \textbf{[N\_success]} / \textbf{[N\_total]}
    \item Failures and reasons: \textbf{[N\_fail + reasons]}
    \item Runtime per document (median/mean): \textbf{[time]}
    \item Output artifact sizes (average): \textbf{[avg JSON sizes]}
\end{itemize}

\subsubsection{Structural baseline results}
\label{sec:eval_baseline_results}

\paragraph{To fill after execution.}
\begin{itemize}
    \item Total missing required sections across corpus: \textbf{[X]}
    \item Most frequently missing sections: \textbf{[Top 5]}
    \item Most frequent extra sections: \textbf{[Top 5]}
\end{itemize}

\subsubsection{Cosine resolver contribution}
\label{sec:eval_cosine_results}

\paragraph{To fill after execution.}
\begin{itemize}
    \item Total rescued mappings: \textbf{[X]}
    \item Reduction in missing-section count: \textbf{[before $\rightarrow$ after]}
    \item Representative rescued mappings: \textbf{[examples]}
\end{itemize}

\subsubsection{Pass~1 results (integrity + semantic adequacy)}
\label{sec:eval_pass1_results}

\paragraph{To fill after execution.}
\begin{itemize}
    \item Total integrity issues detected: \textbf{[X]}
    \item Total semantic issues detected: \textbf{[X]}
    \item Section score distribution summary: \textbf{[histogram summary]}
\end{itemize}

\subsubsection{Pass~2 results (cross-view consistency)}
\label{sec:eval_pass2_results}

\paragraph{To fill after execution.}
\begin{itemize}
    \item Entities extracted totals: \textbf{[counts per type]}
    \item Consistency issues total: \textbf{[X]}
    \item Breakdown by rule category: \textbf{[counts]}
\end{itemize}

\subsection{Comparison of stages}
\label{sec:eval_stage_comparison}

This section synthesizes stage contribution:
\begin{itemize}
    \item what baseline detects and where it over-reports due to naming variation,
    \item what cosine rescue reduces (false ``missing''), and what it additionally detects (placeholders, template copies),
    \item what structure+similarity still misses that Pass~1 detects (semantic inadequacy),
    \item what Pass~1 still misses that Pass~2 detects (cross-view inconsistencies).
\end{itemize}

\subsection{Qualitative analysis and expert feedback}
\label{sec:eval_qualitative}

A qualitative inspection complements quantitative metrics by assessing interpretability and industrial usefulness.

\paragraph{Proposed protocol (to execute if feasible).}
\begin{itemize}
    \item Sample \textbf{[k]} documents covering low/mid/high scores.
    \item For each document, inspect the top \textbf{[m]} findings by severity.
    \item Classify each finding as:
    \begin{itemize}
        \item valid and useful,
        \item valid but low priority,
        \item false positive / ambiguous,
        \item unclear evidence.
    \end{itemize}
\end{itemize}

If feasible, practitioner feedback from a quality manager or architect is incorporated:
\begin{itemize}
    \item perceived usefulness of evidence-backed findings,
    \item acceptance of the score rubric as a prioritization aid,
    \item suggestions for threshold tuning and rule adjustment.
\end{itemize}

\subsection{Threats to evaluation validity}
\label{sec:eval_threats}

The evaluation is subject to several validity threats (discussed in depth in Chapter~8):
\begin{itemize}
    \item \textbf{Construct validity:} proxy metrics (counts/scores) may not fully represent true documentation quality.
    \item \textbf{Internal validity:} threshold sensitivity and prompt formulation may influence outputs.
    \item \textbf{External validity:} results may depend on the company template variant and domain-specific writing conventions.
    \item \textbf{Reliability:} variability in audit passes is mitigated through bounded inputs, schema constraints, and evidence requirements.
\end{itemize}

\subsection{Summary}
\label{sec:eval_summary}

This chapter evaluates the pipeline in terms of stage contribution, auditability, and industrial relevance. Results are reported quantitatively (issue counts, score distributions, and rule-category breakdowns) and qualitatively (evidence-backed examples and practitioner feedback). Chapter~8 interprets the results and discusses strengths, limitations, and industrial applicability.
